%! Functions — Unit 1, Lesson 1.3 — Warm-Up (Answer Key)
\documentclass[10pt]{article}
\usepackage{algebra2-article}
\usepackage{algebra2-key}   % pulls in -boxes; do NOT also load -boxes
\newcommand{\TallMath}[1]{$\displaystyle #1\rule[-1.4em]{0pt}{3.2em}$}

% Pre-drawn coordinate grid; #1 receives the plotted curve and any overlay.
\newcommand{\gridplot}[1]{%
  \begin{tikzpicture}[x=0.36cm, y=0.24cm]
    \draw[step=1, linegray!45, very thin] (-5,-7) grid (5,7);
    \draw[->, charcoal, thick] (-5.4,0) -- (5.4,0) node[right, font=\scriptsize] {$x$};
    \draw[->, charcoal, thick] (0,-7.4) -- (0,7.4) node[above, font=\scriptsize] {$y$};
    \foreach \t in {-4,-2,2,4} {\node[below right=-2pt, font=\tiny] at (\t,0) {$\t$};}
    \foreach \t in {-6,-4,-2,2,4,6} {\node[above left=-2pt, font=\tiny] at (0,\t) {$\t$};}
    #1
  \end{tikzpicture}}

\begin{document}

\pageheader{Unit 1, Lesson 1.3}{Warm-Up --- Answer Key}
\namedateperiod

\vspace{0.06in}

\begin{notesbox}{Getting Ready: Three Ways to Ask One Question}
\small Nothing here is new. Work quickly --- the point is what the three items have in common.

\vspace{3pt}
\textbf{1. Solve (Lesson 1.2).} Find \emph{every} value that makes each expression equal zero.
\begin{center}
\renewcommand{\arraystretch}{1.6}
\begin{tabular}{@{}l l@{}}
(a) \; $2x-6=0$, \quad $x=$ \ans{$3$} &
(b) \; $x^2-2x-3=0$, \quad $x=$ \ans{$3$ or $-1$} \\
\end{tabular}
\end{center}

\vspace{2pt}
\textbf{2. Evaluate (Lesson 1.1).} Each rule below is a recipe. Run it at the two inputs given.
\begin{center}
\renewcommand{\arraystretch}{1.45}
\begin{tabularx}{\linewidth}{@{}X c c@{}}
\hline
\textbf{Rule} & \textbf{Value at $x=0$} & \textbf{Value at $x=3$} \\
\hline
$2x-6$ & \ans{$-6$} & \ans{$0$} \\
$x^2-2x-3$ & \ans{$-3$} & \ans{$0$} \\
$2^x$ & \ans{$1$} & \ans{$8$} \\
\hline
\end{tabularx}
\end{center}
Which rules gave you \textbf{zero} at $x=3$? \ans{the first two} \\
Now look back at item 1. Is that a coincidence? \ans{No} \\
Is there \emph{any} input that makes $2^x$ equal zero? \ans{No} \; (Try a few. Try a negative one.)

\vspace{3pt}
\textbf{3. Read the picture.} Here is a graph. Answer only from the picture.
\begin{center}
\begin{tabular}{@{}c l@{}}
\gridplot{\draw[forest, very thick, domain=-0.6:5, smooth, variable=\x] plot ({\x},{2*\x-6});
          \fill[keyred] (3,0) circle (2.6pt);
          \fill[keyred] (0,-6) circle (2.6pt);}
&
\begin{minipage}[c]{0.52\linewidth}
It crosses the $x$-axis at $x=$ \ans{$3$}
\\[6pt]
It crosses the $y$-axis at $y=$ \ans{$-6$}
\\[6pt]
Left to right, it goes \ans{up} (up / down).
\\[6pt]
Every time $x$ increases by $1$, $y$ changes by \ans{$+2$}
\end{minipage}
\\
\end{tabular}
\end{center}

\vspace{1pt}
\textbf{Put it together.} Item 1(a), the first row of item 2, and the graph above are all the
\emph{same rule}. How can you tell? \ans{Every answer matches: input $3$ gives $0$, input $0$ gives $-6$.}

\end{notesbox}

\begin{teachernote}
\small
\textbf{This warm-up is the lesson in miniature.} The same rule $2x-6$ appears three times --- solved
algebraically, evaluated numerically, and read graphically --- and the answers agree every time. That
agreement \emph{is} \textbf{A.F.1h}: the representations are not three topics, they are one object
seen from three sides. Take the closing question out loud before starting the notes.

\textbf{Item 1 is Lesson 1.2 being cashed in, and today it gets a name.} The values students just
found are about to be called the \textbf{zeros} of a function. Say the definition once here so it
lands as a renaming rather than a new procedure: a zero is a solution of $f(x)=0$, which is the
same thing as an $x$-intercept of the graph.

\textbf{The $2^x$ row is a deliberate plant.} Students will try $x=0$, then negatives, and find
$2^{-1}=\frac12$, $2^{-3}=\frac18$ --- shrinking but never arriving. Do \emph{not} name the
asymptote (that debuts in Unit 5); just get the observation on the board. Section 4 of the notes
needs it, and it is the cleanest contrast between the exponential family and the other two, both of
which hit zero somewhere.

\textbf{Item 3 also pre-loads slope.} ``Every time $x$ increases by $1$, $y$ changes by $+2$'' is
the rate of change, arrived at by counting boxes rather than by a formula --- which is where the
notes pick it up.
\end{teachernote}

\end{document}
